% !TEX root = ../main.tex
\chapter{Calculus of Constructions}

In Chapter 1, we established a pipeline for autonomous theorem proving: 
take a theory, translate its syntax into a Contextual Category, and reframe proof-finding
as the geometric problem of finding a section in a graph.

In this chapter, we scale up the complexity to the Calculus of Constructions (CoC).
Introduced by Coquand and Huet \cite{coquand1988}, CoC is a powerful high-order typed lambda calculus.
It serves as the foundation for the Calculus of Inductive Constructions, the type 
theory powering modern proof assistants like Rocq (formerly Coq).

The primary additions in CoC, compared to Generalized Algebraic Theories (GATs), are: 
\begin{enumerate}
    \item \textbf{Product Types ($\prod$-types):} These are first-class citizens,
        meaning dependent function types exist at the base level rather than just
        in the theory's signature.
    \item \textbf{The Universe of Propositions ($\Prop$):} A type universe whose terms are themselves
        types, allowing for higher-order logic and universal quantification over types.
\end{enumerate}

We define syntax of CoC, map it to a richer categorical structure called
Doctrine of Constructions, and extend our Expansion Algorithm to navigate this new, 
highly complex search space.


\section{Syntactic Formulation}

In this section, we formally define the syntax of Calculus of Constructions.
We use the approach of Streicher in \cite{streicher}, since the original 
definition in \cite{coquand1988} is characterized by its syntax overloading, 
making it much harder to follow.
They just use $[x: A]B$ to present product types, functional abstraction and 
universal quantification. Also, properties are presented as terms of type $\Prop$
and as types themselves. To solve the latter, we will be using a type constructor
$\Proof$ that, given a proposition $p : \Prop$, yields $\Proof(p)$ the type of 
proofs of proposition $p$.

First, we define the class of \textit{pre-well-formed type expressions},
ranged over by capital letters, 
and the class of \textit{pre-well-formed object expressions},
ranged over by lower case letters, by mutual recursion in a BMF-like way: % TODO: add definition of BMF 
\[
\begin{array}{lll}
  E \Coloneqq & (\Pi x : E)\, E & \text{product of types} \\
        & \mid\ \Prop & \text{type of propositions} \\
        & \mid\ \Proof(t) & \text{type of proofs of proposition } t
\end{array}
\]
\[
\begin{array}{lll}
    t \Coloneqq & x & \text{variable} \\
        & \mid\ (\lambda x \colon E) t & \text{functional abstraction} \\
        & \mid\ \App([x \colon E ] E, t, t) & \text{function application} \\
        & \mid\ (\forall x \colon E) t & \text{universal quantification} 

\end{array}
\]
We say a \textit{pre-context} is a syntactical expression of the form 
$x_1 : A_1, \dots, x_n : A_n$, where $x_1, \dots, x_n$ are pairwise 
syntactically different variables. 
Contexts are ranged over by capital greek letters.
We define $\Var(\Gamma)$ as the set of variables declared in $\Gamma$.
We will not distinguish expressions which are equivalent up to renaming of variables
($\alpha$-conversion), to do this we will be using \textit{deBruijn indices}
They are a method to refer to declarations not by name (variables)
but by numbers which tell how far one has to go back into the current context of an expression
to find the declaration of the thing it refers to.
We define a \textit{pre-judgment} to be a syntactic expression of one of the following forms:
\[
\begin{array}{ll}
    A \type & A \text{is a type} \\
    A = B & A \text{ and } B \text{ are equal types} \\
    t : A & t \text{ is an object of type } A \\
    t = s : A & t \text{ and } s \text{ are equal objects of type } $A$
\end{array}
\]
we will use the letter $\mathcal{J}$ to refer to pre-judgments.
We define a \textit{pre-sequent} as a syntactic expression of the following form:
\[
\begin{array}{ll}
    \Gamma \ok & \Gamma \text{ is a context} \\
    \Gamma = \Delta & \Gamma \text{ and } \Delta \text{ are equal contexts} \\
    \Gamma \vdash \mathcal{J} & \text{the judgement } \mathcal{J} \text{ holds w.r.t. } \Gamma
\end{array}
\]


Finally, we define the set of \textit{sequents}, i.e., valid pre-sequents, 
to be generated by the set of rules in \ref{app:derivation-rules-coc}. 
Sequents are ranged over by capital italic letters.
This set of rules is more extent that the original in order to make explicit rules for 
handling contexts and to add more rules stating equality of types and objects.
We can distinguish different kinds of rules for different purposes:
\begin{enumerate}
    \item Context Rules: their purpose is to make contexts first-class citizens of CoC, unlike in GATs where they were defined outside the rule system.
        Theses includes: context formation rules and context equality rules
    \item Structural Rules for Judgments: their purpose is to form and serve as equality rules for judgments.
    \item Rules for Product Types: these rules include everything needed for working with products.
    \item Rules for Propositions: finally, the machinery to work with the $\Prop$ universe, $\Proof$ and $\forall$.
\end{enumerate}

\section{Doctrines of Constructions}

We leverage the concept of context categories introduced in the earlier chapter
to serve as a basic structure for handling dependent types, and add more conditions
so that it can handle products of families of types, universal quantification and the type of properties, in
order to develop a categorical structure that captures the interpretation of the
Calculus of Constructions.

Although products might be added directly in the GAT structure, we enforce them to the base level as
they are needed to the definition of universal quantification.

\begin{definition}
    A \textit{contextual category with products of families of types} is a category $\mathcal{C}$
    together with two mappings $\prod$ and $\eval$ defined over the collection of objects $B$ of 
    $\mathcal{C}$ with $level(B) \geq 2$, such that:
    \begin{enumerate}
        \item Whenever $I \lhd A \lhd B$, then $I \lhd \prod(B)$ and $eval_B \colon p(A)^*\prod(B) \to B$
              is a morphism over $A$
              \[
              \begin{tikzcd}
                p(A)^*\prod(B) \arrow[rr, "\eval_B"] \arrow[dr, "p(p(A)^*\prod(B))"', open]&& B \arrow[dl, "p(B)", open] \\
                & A &
              \end{tikzcd}
              \]
              such that for all $s \in \Sections(B)$ there is a unique section $t \in \Sections(\prod (B))$
              with $\eval_B \circ p(A)^*t = s$. This unique section $t$ is called $\curry(s)$.
              \[
                \begin{tikzcd}
                A \arrow[rr, open] \arrow[ddd, "\id_A"', bend right=85] \arrow[d, "p(A)^*t"] \arrow[dd, "s"', bend right=80] &  & I \arrow[d, "t"'] \arrow[ddd, "\id_I", bend left=80]   \\
                p(A)^*\prod(B) \arrow[d, "\eval_B"] \arrow[rr, "{q(p(A), \prod(B))}"']                                       &  & \prod(B) \arrow[dd, open]                              \\
                B \arrow[d, open]                                                                                                  &  &                                                  \\
                A \arrow[rr, open]                                                                                                 &  & I                                               
                \end{tikzcd}
              \]
              \begin{observation}
                as any section $t$ of $\prod(B)$ gives a section $eval_B \circ p(A)^*t$ of $B$, hence we have a
                canonical 1-1-correspondence of $\Sections(B)$ and $\Sections(\prod(B))$. 
              \end{observation}
        \item For any object $J$ and any morphism $f : J \to I$, the following conditions are satisfied:
              \[
                f^*(\prod(B)) = \prod(f^*B) 
                \quad \text{and} \quad 
                f^*(eval_B) = eval_{f^*B}.
              \]
    \end{enumerate}
\end{definition}

It follows from the definition that substitution preserves functional abstraction or currying:
\begin{lemma}
    Let $\mathcal{C}$ be a contextual category with products of families of types.
    Let $I \lhd A \lhd B$ be objects and $f \colon J \to I$ a morphism in $\mathcal{C}$.
    If $t$ is a section of $B$, then 
    \[
        f^*(\curry(B)) = \curry(f^*B).
    \]
\end{lemma}

Notice that in any contextual category with products of families of types, the following equations hold:
\begin{align*}
    \qquad & \mathrm{eval}_B \circ p(A)^*\mathrm{curry}(s) = s \tag{$\beta$} \\[1em]
    \qquad & \mathrm{curry}(\mathrm{eval}_B \circ p(A)^*t) = t \tag{$\eta$}
\end{align*}
The first equation corresponds to the $\beta$-rule of typed $\lambda$-calculi 
and the second equation corresponds to the $\eta$-rule of typed $\lambda$-calculi.
Contextual Categories with products correspond to $\lambda$-calculus with dependent types.

\begin{definition}
    Let $\mathcal{C}, \mathcal{C}'$ be two contextual categories with products of families of types,
    we a contextual functor $F \colon \mathcal{C} \to \mathcal{C}'$ \textit{preserves products of families
    of types} if for $1 \lhd A \lhd B$ in $\mathcal{C}$, \[
        F\prod(B) = \prod(F B) \quad \text{ and } \quad F \eval_B = \eval_{F B}
    \]
\end{definition}

Now we add the type of propositions. To do it, we find objects in our category whose sections, 
i.e., its terms, correspond to objects.
\begin{definition}
    Let $\mathcal{C}$ be a context category. An \textit{enumeration of types of} $\mathcal{C}$
    is an object $T$ of $\mathcal{C}$ with $\level(T)=2$.

    An enumeration of types $T$ of $\mathcal{C}$ is called a \textit{collection of types} of $\mathcal{C}$
    if for all objects $A$ and morphisms $f, g \colon A \to \father(T)$,
    \[
        f^*T = g^*T \quad \Longrightarrow \quad f = g.
    \]
\end{definition}
It is the object $\father(T)$ that we are interested in.

And finally, by wrapping everything up and adding universal quantification we have the desired structure:
\begin{definition}
    A \textit{doctrine of constructions} is a contextual category with products of families of types
    together with a collection of types $T$ such that for all $A, B \in \Ob \mathcal{C}$ with 
    $A \lhd B$ and $f \colon B \to \Prop$, where $\Prop = \father T$, 
    there exists a necessarily unique morphism $g \colon A \to \Prop$ also denoted by $\forall(f)$
    such that
    \[
        g^*T = \forall(f)^*T = \prod(f^*T).
    \]
\end{definition}

\begin{definition}
    Let $\mathcal{C}, \mathcal{C}'$ be two doctrines of constructions.
    A \textit{functor of doctrines of constructions} is a contextual functor $F \colon \mathcal{C} \to \mathcal{C}'$
    that preserves products of families of types such that if $T$ and $T'$ are the selected collections of types 
    of $\mathcal{C}$ and $\mathcal{C}'$ respectively, $F T = T'$, and for any $A \lhd B$ in $\mathcal{C}$, 
    and $f \colon B \to \Prop$, $F(\forall(f)) = \forall(F(f))$.
\end{definition}

We consider the category $\mathbf{DoC}$ of doctrines of constructions and functors of doctrines of constructions.

Now we see that doctrines of constructions correspond to models of Calculus of Constructions.

\section{Interpretation of Calculus of Constructions}

In this section we define the interpretation of Calculus of Constructions in arbitrary
doctrines of constructions. 

To do this, fixed one doctrine of constructions, an \textbf{interpretation function} $\mathcal{M}$ should be a
mapping such that:
\begin{enumerate}
    \item for every pre-context $\Gamma$ of length $n$, if $\mathcal{M}[\Gamma]$ is defined, then 
          $\mathcal{M}[\Gamma]$ is an object of $\level$ $n$;
    \item for every pre-context of length $n$ and type expression $A$, 
          if $\mathcal{M}[\Gamma ; A]$ is defined then $\mathcal{M}[\Gamma ; A]$ is an object 
          of $\level$ $n + 1$ and $\mathcal{M}[\Gamma] \lhd \mathcal{M}[\Gamma; A]$; 
    \item for every pre-context $\Gamma$ of length $n$ and object expression $t$, 
          if $\mathcal{M}[\Gamma; t]$ is defined then it is a morphism $s : \mathcal{M}[\Gamma] \to A$
          where $A$ is an object of $\level$ $n + 1$ such that \[
          \mathcal{M}[\Gamma] \lhd A \quad \text{and} \quad p(A) \circ s = \id_{\mathcal{M}[\Gamma]}.
          \]
\end{enumerate} 

We want to build this function by induction. We need an auxiliary definition:
\begin{definition}
    By structural induction of expressions we define a function \textbf{$\depth$} mapping
    pre-well-formed expressions and pre-contexts to $\omega$:
    \[\depth[x] = 1\]
    \[\depth[\Prop] = 1\]
    \[\depth[\Proof(A)] = \depth[A] + 1\]
    \[\depth[(\prod x \colon A) B] = \depth[A] + \depth[B] + 1\]
    \[\depth[(\lambda x \colon A) t] = \depth[A] + \depth[t] + 1\]
    \[\depth[(\forall x \colon A) t] =  \depth[A] + \depth[t] + 1\]
    \[\depth[\App([x \colon A] B, t, s)] = \depth[A] + \depth[B] + \depth[t] + \depth[s] + 1\]
    \[
        \text{For } \Gamma \equiv x_1 : A_1, \dots, x_n : A_n, 
        \quad
        \depth[\Gamma] = \sum_{i=1}^{n}\depth[A_i]
    \]
\end{definition}
We define $\mathcal{M}$ by induction on the level of its arguments, where
\[    \level[\Gamma] = \depth[\Gamma], \quad \]
\[    \level[\Gamma; A] = \depth[\Gamma] + \depth[A], \quad \]
\[    \level[\Gamma; t] = \depth[\Gamma] + \depth[t] \]
by the following clauses:
\begin{enumerate}
    \item $\mathcal{M}[\space] = 1$;
    \item $\mathcal{M}[\Gamma, x \colon A] = \mathcal{M}[\Gamma; A]$;
    \item $\mathcal{M}[\Gamma, (\prod x \colon A) B] = \prod(\mathcal{M}[\Gamma, x \colon A; B])$
          if $\mathcal{M}[\Gamma, x \colon A ; B]$ is defined;
    \item $\mathcal{M}[\Gamma, \Prop] = p(\mathcal{M}, 1)^*\Prop$
          if $\mathcal{M}[\Gamma]$ is defined;
    \item $\mathcal{M}[\Gamma; \Proof(p)] = \mathcal{M}[\Gamma; p]^*p(\mathcal{M}[\Gamma], 1)^* T$
          if $\mathcal{M}[\Gamma; p]$ is defined and a section of $p(\mathcal{M}[\Gamma], 1)^*\Prop$
          \[
            \begin{tikzcd}[column sep=3cm]
                \mathcal{M}[\Gamma] 
                    \arrow[dr, "\mathcal{M}{[}\Gamma{;} p{]}"] 
                    \arrow[drr, "q(p(\mathcal{M}{[}\Gamma{]}{,} 1){,} \Prop) \circ \mathcal{M}{[}\Gamma{;} p{]}", bend left=10]
                    \arrow[ddr, "\id_{\mathcal{M}[\Gamma]}", bend right=10]
                    & & \\
                & p(\mathcal{M}[\Gamma], 1)^*\Prop \arrow[r, "q(p(\mathcal{M}{[}\Gamma{]}{,} 1){,}\Prop)"] \arrow[d, open]^*\Prop & \Prop \arrow[d, open]\\
                & \mathcal{M}[\Gamma] \arrow[r, "p(\mathcal{M}{[}\Gamma{]}{,} 1)"] & 1
            \end{tikzcd}
        \]
    \item $\mathcal{M}[\Gamma, x] = s$ 
          provided that $\Gamma \equiv x_1 \colon A_1, \dots, x_n \colon A_n$ and $x_i \equiv x$ 
          and $\mathcal{M}[\Gamma]$ is defined and $\mathcal{M}[\Gamma] = A_n \rhd \dots \rhd A_1 \rhd A_0 \rhd 1$
          and $s$ as follows
          \[
            \begin{tikzcd}[column sep=3cm]
                A_n
                    \arrow[dr, "s"] 
                    \arrow[drr, "{p(A_n, A_i)}", bend left=10]
                    \arrow[ddr, "\id_{A_n}", bend right=10]
                    & & \\
                & p(A_m, A_{i-1})^*A_i \arrow[r, "{q(p(A_n , A_{i-1}), A_i)}"] \arrow[d, open] & A_i \arrow[d, open]\\
                & A_n \arrow[r, "{p(A_n , A_{i-1})}"] & A_{i-1}
            \end{tikzcd}
        \]
        or, in other terms, $\mathcal{M}[\Gamma; x_i] = p(A_n, A_i)^*\Delta(A_i)$ where 
        $\Delta(A_i)$ is the unique arrow $m: A_i \to p(A_i)^*A_i$ such that
        \[
            p(p(A_i)^*A_i) \circ m = q(p(A_i), A_i) \circ m = id_{A_i};
        \]
    \item $\mathcal{M}[\Gamma; (\lambda x \colon A) t] = \curry(\mathcal{M}[\Gamma, x \colon A ; t])$
          if $\mathcal{M}[\Gamma, x \colon A; t]$ is defined.
          By induction hypothesis we have that $\mathcal{M}[\Gamma] \lhd \mathcal{M}[\Gamma; A] \lhd \text{cod}(\mathcal{M}[\gamma, x \colon A ; t])$
          and $s$ is a section of $\mathcal{M}[\Gamma, x \colon A; t]$ and therefore we have:
          \[
                \eval_{\text{cod}(\mathcal{M}[\Gamma, x \colon A; t])} \circ p(\mathcal{M}[\Gamma; A])^*\mathcal{M}[\Gamma; (\lambda x \colon A) t] 
                = \mathcal{M}[\Gamma, x \colon A; t];
          \]
    \item \(
            \mathcal{M}[\Gamma; \App([x \colon A] B), t, s] 
            = 
            \mathcal{M}[\Gamma; s]^*(
                \eval_{\mathcal{M}[\Gamma, x \colon A; B]} 
                \circ p(\mathcal{M}[\Gamma; A])^*
                \mathcal{M}[\Gamma; t]
            )
          \) \\
          provided that $\mathcal{M}[\Gamma ; t]$, $\mathcal{M}[\Gamma; s]$, $\mathcal{M}[\Gamma, x \colon A; B]$
          are defined, and $\mathcal{M}[\Gamma; s]$ is a section of $\mathcal{M}[\Gamma; A]$
          and $\mathcal{M}[\Gamma; t]$ is a section of $\prod(\mathcal{M}[\Gamma, x \colon A; B])$
          \[
          \begin{tikzcd}[column sep=3cm]
            \mathcal{M}[\Gamma] 
                \arrow[r, "{\mathcal{M}[\Gamma; s]}"] 
                \arrow[d, "\substack{
                    \mathcal{M}[\Gamma;s]^*(\eval_{\mathcal{M}[\Gamma, x \colon A; B]} \circ \\
                    p(\mathcal{M}[\Gamma; A])^*\mathcal{M}[\Gamma; t])}"']
            & \mathcal{M}[\Gamma; A] 
                \arrow[d, "{\eval_{\mathcal{M}[\Gamma, x \colon A; B]} \circ p(\mathcal{M}[\Gamma; A])^*\mathcal{M}[\Gamma; t]}"]
            \\
            \mathcal{M}[\Gamma; s]^*\mathcal{M}[\Gamma, x \colon A; B] 
                \arrow[r, "{q(\mathcal{M}[\Gamma; s], \mathcal{M}[\Gamma, x \colon A; B])}"]
                \arrow[d, open]
            & \mathcal{M}[\Gamma, x \colon A ; B] 
                \arrow[d, open]
            \\
            \mathcal{M}[\Gamma]
                \arrow[r, "{\mathcal{M}[\Gamma; s]}"]
            & \mathcal{M}[\Gamma; A]
          \end{tikzcd}
          \]
    \item $\mathcal{M}[\Gamma; (\forall x \colon A) p] = s$
          provided $\mathcal{M}[\Gamma, x \colon A ; p]$ is a section of $p(\mathcal{M}[\Gamma; A], 1)^*\Prop$
          and $s$ is the unique section of $p(\mathcal{M}[\Gamma], 1)$ such that 
          \begin{multline*}
            \forall(q(p(\mathcal{M}[\Gamma; A], 1), \Prop) \circ \mathcal{M}[\Gamma, x \colon A; p])
            = \\
            q(p(\mathcal{M}[\Gamma; A], 1), \Prop) \circ \mathcal{M}[\Gamma, x \colon A; p]
          \end{multline*}
          \[
            \begin{tikzcd}[column sep=3cm]
                \mathcal{M}[\Gamma; A]
                    \arrow[dr, "{\mathcal{M}[\Gamma, x \colon A; p]}"] 
                    \arrow[drr, "{q(p(\mathcal{M}[\Gamma; A], 1), \Prop) \circ \mathcal{M}[\Gamma, x \colon A; p]}", bend left=10]
                    \arrow[ddr, "\id_{\mathcal{M}[\Gamma; A]}", bend right=10]
                    & & \\
                & p(\mathcal{M}[\Gamma; A], 1)^*\Prop \arrow[r, "{q(p(\mathcal{M}[\Gamma; A], 1), \Prop)}"] \arrow[d, open] & \Prop \arrow[d, open]\\
                & \mathcal{M}[\Gamma; A] \arrow[r, "{p(\mathcal{M}[\Gamma; A], 1)}"] & 1
            \end{tikzcd}
        \]
\end{enumerate}

\begin{theorem}[Correctness Theorem]
    For any doctrine of constructions the following correctness conditions, for $\mathcal{M}$ are satisfied:
    \begin{enumerate}
        \item if $\Gamma \ok$ is derivable, then $\mathcal{M}[\Gamma]$ is defined;
        \item if $\Gamma \vdash A \type$ is derivable, then $\mathcal{M}[\Gamma; A]$ is defined;
        \item if $\Gamma \vdash t \colon A$ is derivable, then $\mathcal{M}[\Gamma; A]$ and $\mathcal{M}[\Gamma; t]$
              are both defined and $\mathcal{M}[\Gamma; t]$ is a section of $\mathcal{M}[\Gamma; A]$;
        \item if $\Gamma = \Delta$ is derivable, then $\mathcal{M}[\Gamma]$ and $\mathcal{M}[\Delta]$ are defined 
              and $\mathcal{M}[\Gamma] = \mathcal{M}[\Delta]$;
        \item if $\Gamma \vdash A = B$ is derivable, then $\mathcal{M}[\Gamma; A]$ and $\mathcal{M}[\Gamma; B]$
              are both defined and $\mathcal{M}[\Gamma; A] = \mathcal{M}[\Gamma; B]$;
        \item if $\Gamma \vdash t = s \in A$ is derivable, then $\mathcal{M}[\Gamma; A]$, $\mathcal{M}[\Gamma; t]$
              and $\mathcal{M}[\Gamma; s]$ are defined, both $\mathcal{M}[\Gamma; t]$ and $\mathcal{M}[\Gamma; s]$
              are sections of $\mathcal{M}[\Gamma; A]$ and $\mathcal{M}[\Gamma; t] = \mathcal{M}[\Gamma; s]$.
    \end{enumerate}
\end{theorem}

Hence, doctrines of constructions are models of Calculus of Constructions.

\section{The Term Model of Calculus of Constructions}

In this section we construct a doctrine of constructions $\mathcal{C}_I$ as a \textit{term model}
of calculus of constructions such that the interpretation in this model interprets provably well-defined
contexts, types and objects, and identifies only those contexts, types or objects which are provably equal
up to renaming of the variables bound in contexts.

In more technical terms, this means that for the interpretation function $\mathcal{M}$ assigning meaning
to calculus of constructions in the doctrine of constructions $\mathcal{C}_I$, the following requirements
must be fulfilled:
\begin{enumerate}
    \item for all pre-contexts $\Gamma$ and $\Delta$, if $\mathcal{M}[\Gamma] = \mathcal{M}[\Delta]$,
          then $\Gamma = \Delta$ is derivable;
    \item for all pre-contexts $\Gamma$ and $\Delta$ and pre-type-expressions $A$ and $B$, if 
          $\mathcal{M}[\Gamma; A] = \mathcal{M}[\Delta; B]$, then $\Gamma = \Delta$ and 
          $\Gamma \vdash A = B$ are derivable;
    \item for all pre-contexts $\Gamma$ and $\Delta$, pre-type-expressions $A$ and $B$ and pre-object-expressions
          $t$ and $s$, if $\mathcal{M}[\Gamma; t] = \mathcal{M}[\Delta; s]$ is a section of $\mathcal{M}[\Gamma; A] = \mathcal[\Delta; B]$,
          then $\Gamma = \Delta$, $\Gamma \vdash A = B$ and $\Gamma \vdash t = s \colon A$ are derivable.
\end{enumerate}

In order to be able to identify contexts up to $\alpha$-conversion, we consider a fixed numerable set of variables:
$v_1, \dots, v_n, \dots$. We say that a context $\Gamma \equiv x_1 \colon A_1, \dots, x_n \colon A_n$ is in 
\textbf{$\alpha$-normal form} iff $x_i \equiv v_i$ for all $1 \leq i \leq n$.
Moreover, we  say that a pre-context $\Gamma$ is admissible iff it is in $\alpha$-normal form and $\vdash \Gamma \ok$.

Now we  define $\mathcal{C}_I$:
The objects of $\mathcal{C}_I$ are the admissible contexts modulo provable equivalence, i.e., that we will consider 
two contexts $\Gamma$ and $\Delta$ equal iff $\Gamma = \Delta$ is derivable.
By induction over the structure of derivations, one can see that $\Gamma = \Delta$ implies that $\Gamma$ and
$\Delta$ have the same length as pre-contexts, so we can define the $\level$ of an object of $\mathcal{C}_I$ as the length
of any of its representatives. We will write $[\Gamma]$ to represent the class containing $\Gamma$.
If $\Gamma \equiv v_1 \colon A_1, \dots, v_n \colon A_n$ is an admissible context, then we define
\[
    \father[\Gamma] = [ v_1 \colon A_1, \dots, v_{n-1} \colon A_{n-1} ],
\]
that is well defined according to rule \ref{rule:crefl}1.

If $\Gamma \equiv v_1 \colon A_1, \dots, v_n \colon A_n$ and $\Delta \equiv v_1 \colon B_1, \dots, v_m \colon B_m$
are admissible contexts, a morphism from $[\Gamma]$ to $[\Delta]$ is given by a $m$-tuple $\langle t_1, \dots, t_m \rangle$
of terms such that for $i$, $1 \leq i \leq m$, 
\[
    \Gamma \vdash t_i : B_i[v_1 \mid\ t_1, \dots, v_{i-1} \mid\ t_{i-1}]
\] is derivable,
if $\langle s_1, \dots, s_m \rangle$ is another $m$-tuple such that for  $i$, $1 \leq i \leq m$
\[
    \Gamma \vdash s_i : B_i[v_1 \mid\ s_1, \dots, v_{i-1} \mid\ s_{i-1}]
\] is derivable,
then $\langle t_1, \dots, t_m \rangle$ and $\langle s_1, \dots, s_m \rangle$ are considered equal
iff for $i$, $1 \leq i \leq m$
\[
    \Gamma \vdash t_i = s_i : B_i[v_1 \mid\ s_1, \dots, v_{i-1} \mid\ s_{i-1}]
\] is derivable,
according to rules \ref{rule:crefl}1 and \ref{rule:crefl}2 the definition is independent from the choice of representatives.
We write $[\langle t_1, \dots, t_m \rangle] : [\Gamma] \to [\Delta]$ for the class of $m$-tuples equatable to $\langle t_1, \dots, t_m \rangle$.
If we consider $\Theta \equiv v_1 \colon C_1, \dots, v_k \colon C_k$ is another admissible context and 
$[\langle s_1, \dots, s_k\rangle]: [\Delta] \to [\Theta]$, then the composition of $[\langle t_1, \dots, t_m \rangle]$
with $[\langle s_1, \dots, s_k \rangle]$ is defined as
\[
    [\langle s_1, \dots, s_k \rangle] 
    \circ 
    [\langle t_1, \dots, t_m \rangle] 
    = 
    [\langle 
        s_1[v_1 \mid\ t_1, \dots, v_m \mid\ t_m],
        \dots,
        s_k[v_1 \mid\ t_1, \dots, v_m \mid\ t_m]
    \rangle].
\]
If $\Gamma \equiv v_1 \colon A_1, \dots, v_n \colon A_n$ is an admissible context with $n \geq 1$ then the \textit{canonical projection}
map $p([\Gamma])\colon [\Gamma] \to \father [\Gamma]$ is represented by the tuple $\langle v_1, \dots, v_{n-1} \rangle$.
Let $\Gamma \equiv v_1 \colon A_1, \dots, v_n \colon A_n$ and $\Delta \equiv v_1 \colon B_1, \dots, v_m \colon B_m$ 
admissible contexts and $t = [\langle t_1, \dots, t_m \rangle] \colon [\Gamma] \to [\Delta]$ and 
$C$ a pre-well-formed type expression with $\Delta \vdash C \type$, i.e. $\Delta, v_{m+1} \colon C$ derivable is
an admissible context, then 
\[
    t^*[\Delta, v_{m+1} \colon C] = [\Delta, v_{n+1} \colon C[v_1 \mid\ t_1, \dots, v_m \mid\ t_m]]
\]
\[
    q([t], [\Delta, v_{m+1} \colon C]) = [\langle t_1, \dots, t_m, v_{n+1}] \colon t^*[\Delta, v_{m+1} \colon C] \to [\Delta, v_{m+1} \colon C].
\]
If $\Gamma \equiv v_1 \colon A_1, \dots, v_n \colon A_n, v_{n+1} \colon B, v_{n+2} \colon C$ is an admissible context
we define 
\[
    \prod([\Gamma]) = [\Gamma, v_{n+1} \colon (\prod v_{n+1} \colon B) C]
\]
\[
    ev_{[\Gamma]} = [\langle v_1, \dots, v_n, v_{n+1}, \App([v_{n+1} \colon B]C, v_{n+2}, v_{n+1})\rangle].
\]
which can be seen to be a morphism from $[\Gamma, v_{n+1} \colon B, v_{n+2} \colon (\prod v_{n+1} B) C]$ to $[\Gamma]$
over the object $[v_1 \colon A_1, \dots, v_n \colon A_n]$.
The object $\Prop$ of $C_I$ is defined as $[v_1 \colon \Prop]$ and the object $T$ of $C_I$ is defined as the object
$[v_1 \colon \Prop, v_2 \colon \Proof(v_1)]$.
If $\Gamma \equiv v_1 \colon A_1, \dots, v_n \colon A_n, v_{n+1} \colon B$ and $[\langle p \rangle]$ is a morphism 
from $[\Gamma]$ to $[v_1 \colon \Prop]$ then we define
\[
    \forall([p]) = [\langle (\forall v_{n+1} \colon B) p \rangle] 
        \colon [v_1 \colon A_1, \dots, v_n \colon A_n] 
        \to [v_1 \colon \Prop].
\]
It is proven in \cite{streicher} that $C_I$ is in fact a doctrine of constructions. Also, it is immediate to see
that $C_I$ does not interpret or identify more objects or terms that those that are provably defined as equal.

\begin{theorem} [Completeness Theorem]
    Let $\mathcal{M}$ be the interpretation of calculus of construction in $C_I$. Then 
    \begin{enumerate}
        \item if $\mathcal{M}[\Gamma]$ is defined then $\Gamma \ok$ is derivable;
        \item if $\mathcal{M}[\Gamma ; A]$ is defined then $\Gamma \vdash A \type$ is derivable;
        \item if $\mathcal{M}[\Gamma; t]$ is defined then $\Gamma \vdash t \colon A$ is derivable for some pre-type-expression $A$;
        \item if $\Gamma \equiv x_1 \colon A_1, \dots, x_n \colon A_n$, $\Delta \equiv y_1 \colon B_1, \dots, y_m \colon B_m$
              and $\mathcal{M}[\Gamma; A] = \mathcal{M}[\Delta; B]$, then $n = m$ and 
              \begin{multline*} 
                    v_1 \colon A_1 [x_1 \mid\ v_1, \dots, x_n \mid\ v_n], \dots, v_n \colon A_n[x_1 \mid\ v_1, \dots, x_n \mid\ v_n]
                = \\ =
                    v_1 \colon B_1 [y_1 \mid\ v_1, \dots, y_m \mid\ v_m], \dots, v_m \colon B_m[y_1 \mid\ v_1, \dots, y_m \mid\ v_m]
              \end{multline*}
              \begin{multline*}
                    v_1 \colon A_1 [x_1 \mid\ v_1, \dots, x_n \mid\ v_n], \dots, v_n \colon A_n[x_1 \mid\ v_1, \dots, x_n \mid\ v_n] \\
                \vdash
                    A[x_1 \mid\ v_1, \dots, x_n \mid\ v_n] = B[y_1 \mid\ v_1, \dots, y_m \mid\ v_m]
              \end{multline*}
              are derivable;
        \item if $\Gamma \equiv x_1 \colon A_1, \dots, x_n \colon A_n$, $\Delta \equiv y_1 \colon B_1, \dots, y_m \colon B_m$
              and $\mathcal{M}[\Gamma; t] = \mathcal{M}[\Delta, s]$, then $n = m$ and
              \begin{multline*} 
                    v_1 \colon A_1 [x_1 \mid\ v_1, \dots, x_n \mid\ v_n], \dots, v_n \colon A_n[x_1 \mid\ v_1, \dots, x_n \mid\ v_n]
                = \\ =
                    v_1 \colon B_1 [y_1 \mid\ v_1, \dots, y_m \mid\ v_m], \dots, v_m \colon B_m[y_1 \mid\ v_1, \dots, y_m \mid\ v_m]
              \end{multline*} is derivable,
              and, for some pre-type-expression $A$,
              \begin{multline*}
                    v_1 \colon A_1 [x_1 \mid\ v_1, \dots, x_n \mid\ v_n], \dots, v_n \colon A_n[x_1 \mid\ v_1, \dots, x_n \mid\ v_n] \\
                \vdash
                    t[x_1 \mid\ v_1, \dots, x_n \mid\ v_n] = s[y_1 \mid\ v_1, \dots, y_m \mid\ v_m] : A
              \end{multline*}
              is derivable.
    \end{enumerate}
\end{theorem}

It can be also shown that $\mathcal{C}_I$ is the initial object of the category of doctrines of constructions and 
structure preserving functors, named $\mathbf{DoC}$. Hence tha term model satisfies the properties of soundness, 
completeness and initiality that we wanted.

\section{Extending Calculus of Constructions and its \\Interpretation} \label{sect:extending-calculus-of-constructions-and-its-interpretation}

In this section we show how to add symbols to Calculus of Constructions, the same thing that we could do in 
GATs. We need to understand the current CoC as just as a minimal starting point, same as the empty theory for GATs. 
Inspired by the constructions for Categories with Families presented in \cite{bezem2021}, we inductively
define the concept of signature $\Sigma$ (corresponding to theory when talking about GATs), CoC extended by $\Sigma$,
the corresponding category $\mathbf{DoC}_\Sigma$ and interpretation functions $\mathcal{M}^\Sigma$.

As a starting case, we consider the empty signature $\Sigma = \emptyset$, together with $\text{CoC}_\emptyset$
representing the current syntax of Calculus of Constructions. The category $\mathbf{DoC}_\emptyset$ is exactly $\mathbf{DoC}$,
and the notion of interpretation morphism stays the same.

Now, for the inductive step, assume we have defined $\Sigma$ as a valid signature, and the associated syntax
$\text{CoC}_\Sigma$, category $\mathbf{DoC}_\Sigma$ and interpretation function $\mathcal{M}$. Then 
\begin{itemize}
    \item \textbf{Adding a sort symbol:} 
          Let $\Gamma$ be a context in $\text{CoC}_\Sigma$. We can extend $\Sigma$
          with a sort symbol $S$ relative to $\Gamma$, to obtain the signature $\Sigma' = (\Sigma, (\Gamma, S))$.
          Then, we extend $\text{CoC}_\Sigma$ to obtain $\text{CoC}_{\Sigma'}$ by adding a production of 
          pre-type-expressions $E \Coloneqq S$ and the inference rule $\Gamma \vdash S$.
          The category $\mathbf{DoC}_{\Sigma'}$ has as objects pairs $(\mathcal{C}, S_\mathcal{C})$, where 
          $\mathcal{C}$ is an object of $\mathbf{DoC}_\Sigma$ and $\mathcal{M}_\mathcal{C}[\Gamma] \lhd A_\mathcal{C}$.
          Finally, we define $\mathcal{M}'_\mathcal{C}$, for each object $\mathcal{C}$ of $\mathbf{DoC}_{\Sigma'}$,
          by adding $\mathcal{M}'_\mathcal{C}[\Gamma; S] = S_\mathcal{C}$.
    \item \textbf{Adding an operator symbol:}
          Let $\Gamma$ be a context and $A$ be a type under $\Gamma$ of $\text{CoC}_\Sigma$,
          we can extend $\Sigma$ with a new operator symbol $f$ relative to $\Gamma$ and $A$
          obtaining $\Sigma' = (\Sigma, (\Gamma, A, f))$.
          Then, we extend $\text{CoC}_\Sigma$ to obtain $\text{CoC}_{\Sigma'}$ by adding a production of 
          pre-object-expressions $t \Coloneqq f$ and the inference rule $\Gamma \vdash f \colon A$.
          The category $\mathbf{DoC}_{\Sigma'}$ has as objects pairs $(\mathcal{C}, f_\mathcal{C})$,
          where $\mathcal{C}$ is an object of $\mathbf{DoC}_\Sigma$ and $f_\mathcal{C}$ is a section of $\mathcal{M_C}[\Gamma, x \colon A]$.
          Finally, we define $\mathcal{M}'_\mathcal{C}$, for each object $\mathcal{C}$ of $\mathbf{DoC}_{\Sigma'}$,
          by adding $\mathcal{M}'_\mathcal{C}[\Gamma; f] = f_\mathcal{C}$.
    \item \textbf{Adding an equality of types:}
          Let $\Gamma$ be a context, $A, A'$ types under $\Gamma$ of $\text{CoC}_\Sigma$, we can extend $\Sigma$ with
          a new equation $A = A'$ relative to $\Gamma$, to obtain $\Sigma' = (\Sigma, (\Gamma, A, A'))$. % TODO: is this good notation?
          Then, we extend $\text{CoC}_\Sigma$ to obtain $\text{CoC}_{\Sigma'}$ by adding the inference rule 
          $\Gamma \vdash A = A'$.
          The category $\mathbf{DoC}_{\Sigma'}$ is a full subcategory of $\mathbf{DoC}_\Sigma$ containing the objects 
          $\mathcal{C}$ of $\mathbf{DoC}_\Sigma$ such that $\mathcal{M_C}[\Gamma; A] = \mathcal{M_C}[\Gamma; A']$.
          $\mathcal{M}'$ is the restriction of $\mathcal{M}$ to the objects of $\mathbf{DoC}_{\Sigma'}$.
    \item \textbf{Adding an equality of terms:}
          Let $\Gamma$ be a context, $A$ a type under $\Gamma$ and $a, a'$ terms of type $A$ under $\Gamma$ 
          of $\text{CoC}_\Sigma$. We can extend $\Sigma$ with a new equation $a = a'$ relative to $\Gamma$ and $A$,
          to obtain $\Sigma' = (\Sigma, (\Gamma, A, a, a'))$.
          Then, we extend $\text{CoC}_\Sigma$ to obtain $\text{CoC}_{\Sigma'}$ by adding the inference rule 
          $\Gamma \vdash a = a' \colon A$.
          The category $\mathbf{DoC}_{\Sigma'}$ is a full subcategory of $\mathbf{DoC}_\Sigma$ containing the objects 
          $\mathcal{C}$ of $\mathbf{DoC}_\Sigma$ such that $\mathcal{M_C}[\Gamma; a] = \mathcal{M_C}[\Gamma; a']$.
          $\mathcal{M}'$ is the restriction of $\mathcal{M}$ to the objects of $\mathbf{DoC}_{\Sigma'}$.
\end{itemize}

Under this construction, we can consider $\mathcal{C}_\Sigma$ to be the exact same construction as $\mathcal{C}_I$,
but with $\text{CoC}_\Sigma$ instead of $\text{CoC}_\emptyset$. Following a similar proof that the original proof
for $\mathcal{C}_I$ in \cite{streicher} we can prove the corresponding \textit{Completeness Theorems} and 
that $\mathcal{C}_\Sigma$ is the initial object of $\mathbf{DoC}_\Sigma$. 

\section{Redefinition of the Problem}

Now, we want to redefine our problem from the previous definition for GATs to our new syntax, then rephrase it to 
similar categorical terms such that an extended version of the expansion algorithm works.

The problem now can be formalized as follows:
Given a valid signature of Calculus of Constructions $\Sigma$ and a derivable sequent $\Gamma \vdash A \type$ in 
$\text{CoC}_\Sigma$, we want to find a pre-object-expression $t$ such that 
$\Gamma \vdash t \colon A$ is a derivable rule in $\text{CoC}_\Sigma$.

We can consider this from the point of view of the category $\mathcal{C}_\Sigma$ introduced in the previous section.
This means that our problem is equivalent to finding a section $s$ of 
$\mathcal{M}[\Gamma; A] = [\Gamma, v_{n+1} \colon A]$ in $\mathcal{C}_\Sigma$, assuming
$\Gamma$ is a context of length $n$ in normal form.

Moreover, we want to use the reduction of our search space we did on the previous case
by considering the category over a base object. To do it we need the following result:

\begin{definition}
    Given a doctrine of constructions $\mathcal{C}$, and an arbitrary object $A$ of $\mathcal{C}$. The \textit{relativization}
    of $\mathcal{C}$ to $A$, noted as $\mathcal{C}_A$, is defined as follows:
    \begin{enumerate}
        \item $\mathcal{C}_A$ is a category whose objects are the collection of objects $B$ over $A$ in $\mathcal{C}$,
              i.e. objects $B$ in $\mathcal{C}$ such that $A \leq B$;
        \item for two objects $B, C$ in $\mathcal{C}_A$, a morphism from $B$ to $C$ over $A$, i.e., a morphism in $\mathcal{C}_A$,
              is a morphism $f \colon B \to C$ in $\mathcal{C}$ such that $p(C, A) \circ f = p(B, A)$, composition
              of morphisms is inherited from $\mathcal{C}$;
        \item $\father_A(B) = \father(B)$ for $B > A$ and $\father_A(A) = \father(A)$;
        \item for objects $B > A$, we define $p_A(B) = p(B) \colon B \to \father_A(B)$;
        \item for $B, C, D$ objects of $\mathcal{C}$ with $\father_A(D) = C$ and morphism $f \colon B \to C$ in $\mathcal{C}_A$
              we define \[
                f^{*_A}D = f^*D \quad \text{ and } \quad q_A(f, D) = q(f, D);
              \]
        \item for $B$ in $\mathcal{C}$ with $\level_A(B) \geq 2$,
              \[
                \prod_A(B) = \prod(B) \quad \text{ and } \quad \eval^A_B = \eval_B;
              \]
        \item if $T$ is the collection of types of $\mathcal{C}$, then 
              \[
                T_A = p(A)^*T
              \]
              is the collection of types of $\mathcal{C}_A$.
    \end{enumerate}
\end{definition}
\begin{theorem}
    The relativization $\mathcal{C}_A$ of a doctrine of constructions $\mathcal{C}$ to an object $A$ is a doctrine of constructions.
\end{theorem}
\begin{proof}
    We already know from the first chapter that $\mathcal{C}_A$ is a contextual category. Let us check if it is also a 
    doctrine of constructions with these extra structure we introduce.

    It is clear that $\mathcal{C}_A$ is a contextual category with products of families of types as for any object $B$ of $\mathcal{C}_A$ 
    with $\level_A(B) \geq 2$, we have that $\level(B) \geq 2$ and as $\mathcal{C}$ is a contextual category with products of families of 
    types, $\prod_A(B) = \prod(B)$ is defined and $\father(\prod (B)) = \father^2 (B)$, and hence $\prod_A(B) \rhd A$, hence it is well defined.
    Then for $\eval_B^A = \eval_B \colon p(\father(B))^*\prod(B) \to B$ is a morphism in $\mathcal{C}_A$ as by definition the following 
    diagram commutes: \[
        \begin{tikzcd}
            p(\father(B))^*\prod(B) 
                \arrow[rr, "{\eval^A_B = \eval_B}"]
                \arrow[rd, open]
                \arrow[rdd, open]
            & & B 
                \arrow[ld, open]
                \arrow[ldd, open]
            \\
            & \father(B) \arrow[d, open] & \\
            & A &
        \end{tikzcd}
    \]
    Now, if $s$ is a section of $B$ in $\mathcal{C}_A$, it is a section of $B$ in $\mathcal{C}$. Hence there exists a unique section $t$ of $\prod(B)$
    in $\mathcal{C}$ such that $eval_B \circ p(A)^*t = s$, as $p(\prod(B)) \circ t = \id_{\father(\prod(B))}$, it is clear that $t$ is a section of $\prod_A(B)$
    in $\mathcal{C}_A$.
    Now if $J$ is an object and $f \colon J \to \father^2(B)$ is a morphism in $\mathcal{C}_A$, in particular, they belong to $\mathcal{C}$.
    So we have that $\prod(f^*B) = f^*\prod(B)$ and $\eval_{f^*B} = f^*\eval_B$, but by definition of $\prod_A$ and $\eval_A$ we have that
    $\prod_A(f^*B) = f^*\prod_A(B)$ and $\eval^A_{f^*B} = f^*\eval^A_B$. With this we conclude that $\mathcal{C}_A$ is a contextual category with
    products of families of types.

    Now we see that $T_A = p(A)^*T$ is indeed a collection of types of $\mathcal{C}_A$. Consider then two morphisms $f, g \colon J \to \father(T_A)$,
    for some object $A$ in $\mathcal{C}_A$, such that $f^*T_A = g^*T_A$. We want to see that indeed $f = g$. We have that $f^*(q(p(A),\Prop)^*T) = g^*(q(p(A), \Prop)^*T)$
    and we derive that $(q(p(A), \Prop) \circ f)^*T = (q(p(A), \Prop) \circ g)^*T$. As $T$ is a collection of types in $\mathcal{C}$,
    $q(p(A), \Prop) \circ f = q(p(A), \Prop) \circ g$. And finally $f = g$ as both make the following pullback diagram commute: \[
        \begin{tikzcd}[column sep=3cm]
            J 
                \arrow[rd, "{f = g}"]
                \arrow[rrd, "{q(p(A), \Prop) \circ f}", bend left = 10]
                \arrow[rdd, open, bend right = 10] 
            & & \\
            & \Prop_A 
                \arrow[r, "{q(p(A), \Prop)}"]
                \arrow[d, open]
            & \Prop
                \arrow[d, open]
            \\
            & A 
                \arrow[r, open]
            & 1
        \end{tikzcd}
    \]

    Finally, let $B, C$ be objects of $\mathcal{C}_A$ and $f \colon C \to \Prop_A$ a morphism over $A$. Then we want to see that there exists a
    necessarily unique morphism $g \colon B \to \Prop_A$ such that $g^*T_A = \prod_A(f^*T_A)$.
    As $\mathcal{C}$ holds this property, we consider $q(p(A), \Prop) \circ f \colon C \to \Prop$, hence there exists a necessarily unique $g' \colon B \to \Prop$
    such that $g'^*T = \prod ((q(p(A), \Prop) \circ f)^* T) = \prod (f^*T_A)$, and by considering $g$ the unique function that makes the following pullback diagram commute:  \[
        \begin{tikzcd}[column sep=3cm]
            B
                \arrow[rd, "g"]
                \arrow[rrd, "g'", bend left = 10]
                \arrow[rdd, open, bend right = 10] 
            & & \\
            & \Prop_A 
                \arrow[r, "{q(p(A), \Prop)}"]
                \arrow[d, open]
            & \Prop
                \arrow[d, open]
            \\
            & A 
                \arrow[r, open]
            & 1
        \end{tikzcd}
    \] we then have that $g' = q(p(A), \Prop) \circ g$, hence: 
    \begin{align*}
        g^*T_A &= g^*(q(p(A), \Prop)^*T)  \\
               &= (q(p(A), \Prop) \circ g)^*T \\
               &= g'^*T  \\
               &= \prod((q(p(A), \Prop) \circ f)^*T) \\
               &= \prod(f^*(q(p(A), \Prop)^*T)) = \prod_A(f^*T_A)
    \end{align*}
    This concludes the proof as $g$ is a morphism over $A$.
\end{proof}


With this result, and with the notation of $\Gamma$ subscript introduced in the second chapter, we reduce our problem to 
finding a section of $[ \Gamma, x \colon A ]$ in $\mathcal{C}_{\Sigma, \Gamma}$, the relativization of $\mathcal{C}_\Sigma$
to $\Gamma$.

\section{Expansion Algorithm}

Now we extend the expansion algorithm from the previous chapter. We reintroduce the rules in order to match the syntax we are 
using for Calculus of Constructions.
First, consider $G$ to be a finite subcategory of $\mathcal{C}_{\Sigma, \Gamma}$ containing $[\Gamma]$,
$[\Gamma, p \colon \Prop]$, $\Gamma, x \colon A$ and their projections.
Now we define the expansion rules. If $\Delta$ is a context, then:
\begin{enumerate}
    \ruleitem{E1}{sort-symbol-intro-2}
        For each sort symbol in the signature $\Sigma$, with an introductory rule $\Omega \vdash S \type$,
        such that $\Delta \equiv \Gamma, y_1 \colon B_1, \dots, y_m \colon B_m$, $\Omega \equiv z_1 \colon C_1, \dots, z_w \colon C_w$ and
        $G_n$ contains sections
        \begin{align*}
            [ \Delta ] 
                &\xrightarrow{\langle y_1, \dots, y_m, t_1 \rangle_\Gamma} 
                [ \Delta, z_1 \colon C_1 ] \\
            [ \Delta ] 
                &\xrightarrow{\langle y_1, \dots, y_m, t_2 \rangle_\Gamma} 
                [ \Delta, z_2 \colon C_2[z_1 \mid\ t_1] ] \\
            & \vdots \\
            [ \Delta ] 
                &\xrightarrow{\langle y_1, \dots, y_m, t_w \rangle_\Gamma} 
                [ \Delta, z_w \colon C_w[z_1 \mid\ t_1, \dots, z_{w-1} \mid\ t_{w-1}] ].
        \end{align*}
        Then we add $[ \Delta, y \colon A[z_1 \mid\ t_1, \dots, z_w \mid\ t_w] ]$ and its projection to $G_{n+1}$.
    \ruleitem{E2}{variable-duplication-2}
        If $\Delta \equiv \Gamma, y_1 \colon B_1, \dots, y_m \colon B_m$, then we add the sections 
        \[
            [\Delta] \xrightarrow{\langle y_1, \dots, y_m, y_i \rangle_\Gamma} [\Delta, y \colon B_i] 
        \]\[
            [\Delta] \xrightarrow{\langle y_1, \dots, y_m, x_j \rangle_\Gamma} [\Delta, x \colon A_j]
        \]
        and their corresponding projections for each $1 \leq i \leq m$ and each $1 \leq j \leq n$.
    \ruleitem{E3}{operator-symbol-intro-2}
        For each sort symbol in the signature $\Sigma$, with an introductory rule $\Omega \vdash F \colon A$,
        such that $\Delta \equiv \Gamma, y_1 \colon B_1, \dots, y_m \colon B_m$, $\Omega \equiv z_1 \colon C_1, \dots, z_w \colon C_w$
        and the following sections exist in $G_n$
        \begin{align*}
            [ \Delta ]
                &\xrightarrow{\langle y_1, \dots, y_m, t_1 \rangle_\Gamma}
                [ \Delta, z_1 \colon C_1 ] \\
            [ \Delta ]
                &\xrightarrow{\langle y_1, \dots, y_m, t_2 \rangle_\Gamma}
                [ \Delta, z_2 \colon C_2[z_1 \mid\ t_1] ] \\
            & \vdots \\
            [ \Delta ]
                &\xrightarrow{\langle y_1, \dots, y_m, t_w\rangle_\Gamma}
                [ \Delta, z_w \colon C_w[z_1 \mid\ t_1, \dots, z_{w-1} \mid\ t_{w-1}] ]
        \end{align*}
        we add the section
        \[
            [ \Delta ]
            \xrightarrow{\langle y_1, \dots, y_m, F [z_1 \mid\ t_1, \dots, z_w \mid\ t_w] \rangle_\Gamma}
            [ \Delta, z \colon A[z_1 \mid\ t_1, \dots, z_w \mid\ t_w] ].
        \]
    \ruleitem{E4}{pullback-2}
        If $\Delta \equiv \Gamma, y_1 \colon B_1, \dots, y_m \colon B_m$, 
        for every morphism \[
            [\langle t_1, \dots, t_{m-1} \rangle_\Gamma]
            : [\Gamma, z_1 \colon C_1, \dots, z_w \colon C_w ]
            \to [\Gamma y_1 \colon B_1, \dots, y_{m-1} \colon B_{m-1} ]
        \]
        we construct the corresponding pullback:
        \[
        \begin{tikzcd}
            \Gamma, 
            z_1 \colon C_1, \dots, z_w \colon C_w,
            y_m \colon B_m[y_1 \mid\ t_1, \dots, y_{m-1} \mid\ t_{m-1}]
                \arrow[r, "{\langle t_1, \dots, t_{m-1}, y_m \rangle_\Gamma}"]
                \arrow[d, open] 
            & 
            \Gamma,
            y_1 \colon B_1, \dots, y_m \colon B_m
                \arrow[d, open]
            \\ 
            \Gamma,
            z_1 \colon C_1, \dots, z_w \colon C_w 
                \arrow[r, "{\langle t_1, \dots, t_{m-1} \rangle_\Gamma}"]
            &
            \Gamma,
            y_1 \colon B_1, \dots, y_{m-1} \colon B_{m-1}
        \end{tikzcd}
        \]
        to $G_{n+1}$.
    \ruleitem{E5}{product-formation-2}
        If $\Delta \equiv \Gamma, v_1 \colon B_1, \dots, v_m \colon B_m, v_{m+1} \colon C, v_{n+2} \colon D$
        we add \[
            \prod([\Delta]) = [
                \Gamma,
                v_1 \colon B_1, \dots, v_m \colon B_m,
                v_{m+1} \colon (\prod v_{m+1} \colon C) D
            ]
        \]
        \[
            \eval_{[\Delta]} = [
                \Gamma,
                v_1, \dots, v_{m+1},
                \App([v_{m+1} \colon C]D, v_{m+2}, v_{m+1})
            ].
        \]
        a morphism from $[\Gamma, v_1 \colon A_1, \dots, v_m \colon B_m, v_{n+1} \colon C, v_{n+2} \colon (\prod v_{n+1} \colon B)C ]$
        to $[\Delta]$, and the corresponding projections.
    \ruleitem{E6}{proof-forall-intro-2}
        If $\Delta \equiv \Gamma, v_1 \colon B_1, \dots, v_m \colon B_m, v_{m+1} \colon C$,
        for any morphism $[\langle p \rangle_\Gamma] \colon [\Delta] \to [\langle v_1 \colon \Prop \rangle_\Gamma]$
        we add \[ 
            \forall([\langle p \rangle_\Gamma]) 
            = [ \langle (\forall v_{m+1} C)p \rangle_\Gamma]
            \colon [\Gamma, v_1 \colon B_1, \dots, v_m \colon B_m]
            \to [ \langle v_1 \colon \Prop \rangle_\Gamma]
        \]
        \[
            [\Delta, x \colon \Proof(p)]
        \]
        and its corresponding projection. 
\end{enumerate}
After each expansion, we complete the category in the following way:
for every product object, i.e., for each object $[\Delta]$ with $\Delta$
of the form $\Omega, (\prod x \colon A) B$, and every section 
$s = \langle x, a \rangle_\Omega$ of $[\Omega, x \colon A, y \colon B]$
we add $\curry(s) = \langle (\lambda x \colon A) a \rangle_\Omega$
as a section of $[\Delta]$.
Also, we add the composition of morphisms to $G_{n+1}$ in order to complete it as a subcategory of $\mathcal{C}_{\Sigma, \Gamma}$.

\begin{theorem}
    Every object and every morphism of $\mathcal{C}_{\Sigma, \Gamma}$ are constructible by expansion.
\end{theorem}
\begin{proof} % TODO: how do Thin and Sub rules affect?
    We prove this result by induction over the $\level$ of contexts, pre-type-expressions and terms.

    Consider a context $\Delta \equiv \Gamma, y_1 \colon B_1, \dots, y_m \colon B_m$ in $\text{CoC}_\Sigma$, and the corresponding
    object $[\Delta]$ in $\mathcal{C}_{\Sigma, \Gamma}$. We prove that 
    this object is constructible by expansion. 
    We consider how the sequent $\Delta \ok$ was derived. It has to be by one of the following rules: 
    \ref{rule:empty}, \ref{rule:cont-intro}, \ref{rule:cont-thin}, \ref{rule:cont-sub}, \ref{rule:crefl} and \ref{rule:reflection1}.
    Since the case of \ref{rule:crefl} and \ref{rule:reflection1} are trivial, we consider each of the other cases:
    \begin{enumerate}
        \item[\ref{rule:empty}] This is the base case, when $\Delta$ is the empty context.
            This may only happen when $\Gamma = \Delta = \langle \rangle$ and $m = 0$, and as the object $[\Gamma]$
            is included in $G$, it is always constructible by expansion.
        \item[\ref{rule:cont-intro}] This means that the sequent $\Gamma, y_1 \colon B_1, \dots, y_{m-1} \colon B_{m-1} \vdash A_m \type$
            is derivable. 
            Which itself can % TODO: marked as ? by Carles
            only be derived using the following rules: 
            \ref{rule:cont-intro} (in the reverse form), \ref{rule:refl-typ}, \ref{rule:prod-form}, \ref{rule:prop}, \ref{rule:prop-type}
            and any introductory rule of a sort symbol in $\Sigma$.
            The derivation from rules \ref{rule:cont-intro} and \ref{rule:refl-typ} are trivial, as they both enforce the existence of the object $[\Delta]$ by premise,
            hence we will consider the other cases:
            \begin{enumerate}
                \item[\ref{rule:prod-form}] Then, we have that $A_m$ is a pre-type-expression of the form $(\prod x \colon N)M$ 
                    for some variable $x$ and two pre-type-expressions $N, M$ such that the sequent 
                    $\Gamma, y_1 \colon B_1, \dots, y_{m-1} \colon B_{m-1}, x \colon N \vdash M \type$ is derivable.
                    By induction hypothesis, the object 
                    \begin{multline*}
                        \mathcal{M}[
                        \Gamma, 
                        y_1 \colon B_1, \dots, y_{m-1} \colon B_{m-1},
                        x \colon N ; M
                        ] 
                        = \\ =
                        [
                            \Gamma, 
                            y_1 \colon B_1, \dots, y_{m-1} \colon B_{m-1}, 
                            x \colon N, y \colon M
                        ]
                    \end{multline*}
                     is constructible by expansion. By the expansion rule \ref{rule:product-formation-2}
                     applied to that object, gives us that 
                     \begin{multline*}
                        \prod([
                            \Gamma,
                            y_1 \colon B_1, \dots, y_{m-1} \colon B_{m-1},
                            x \colon N, y \colon M
                        ]) 
                        = \\ =
                        [
                            \Gamma,
                            y_1 \colon B_1, \dots, y_{m-1} \colon B_{m-1},
                            x \colon (\prod x \colon N) M
                        ] 
                     \end{multline*}
                     is constructible, that is equal to $[\Delta]$ as they just differ by the naming of bound variables.
                     Hence, $[\Delta]$ is constructible by expansion.
                \item[\ref{rule:prop}] Then we have that $A_m$ is equal to $\Prop$ as a pre-type-expression, and 
                     and, with $\Delta' \equiv \Gamma, y_1 \colon B_1, \dots, y_{m-1} \colon B_{m-1}$, $\Delta' \ok$ is a derivable sequent.
                     By induction hypothesis, we know that $[\Delta']$ 
                     is constructible by expansion. Also, as $G$ contains $[\Gamma, x \colon \Prop]$ and it's projection to 
                     $[\Gamma]$, by applying the expansion rule \ref{rule:pullback-2} to $[\Gamma, x \colon \Prop]$
                     with the projection morphism from $[\Delta']$ to $[\Gamma]$,
                     we get that the following pullback diagram is constructible by expansion:
                     \[
                     \begin{tikzcd}
                        {[\Delta', x \colon \Prop]}
                            \arrow[r, "{\langle x\rangle_\Gamma}"] 
                            \arrow[d, open]
                        & {[\Gamma, x \colon \Prop]}
                            \arrow[d, open]
                        \\
                        {[\Delta']}
                            \arrow[r, open]
                        & {[\Gamma]}
                     \end{tikzcd}
                     \]
                     and $[\Delta', x \colon \Prop] = [\Delta]$ as they only differ by the naming of the variables.
                \item[\ref{rule:prop-type}] Then we have that there is a pre-object-expression $p$ such that $A_m$ equals $\Proof(p)$ as a pre-type-expression
                     and the sequent $\Delta' \vdash p \colon \Prop$, with $\Delta' \equiv \Gamma, y_1 \colon B_1, \dots, y_{m-1} \colon B_{m-1}$, is derivable.
                     By induction hypothesis, $[\Delta']$ is constructible by expansion.
                     Moreover, by applying induction hypothesis,
                     \[
                        \mathcal{M}[\Delta'; p] = \langle p \rangle_\Gamma
                     \]
                     is constructible by expansion.
                     Finally, by applying expansion rule \ref{rule:proof-forall-intro-2},
                     we get that $[\Delta] = [\Delta', x \colon \Proof(p)]$ is constructible by expansion.
                \item[\textsc{Introductory Rule}] Consider a sort symbol $S$ with the corresponding introductory rule $\Omega \vdash S$,
                     with $\Omega \equiv z_1 \colon C_1, \dots, z_k \colon C_k$, and $t_1, \dots, t_k$ pre-object-symbols
                     such that $A_m = S[z_1 \mid\ t_1, \dots, z_k \mid\ t_k]$.
                     Then, for each $i$, $1 \leq i \leq k$, $\Delta' \vdash t_i \colon C_i [z_1 \mid\ t_1, \dots, z_{i-1} \mid\ t_{i-1}]$
                     are derivable judgments, where $\Delta' \equiv \Gamma, y_1 \colon B_1, \dots, y_{m-1} \colon B_{m-1}$.
                     Hence, by inductive hypothesis, the sections \[
                        \mathcal{M}[\Delta'; t_1] \colon [\Delta'] \to [\Delta', z_1 \colon C_1]
                     \] \[
                        \mathcal{M}[\Delta'; t_2] \colon [\Delta'] \to [\Delta', z_2 \colon C_2[z_1 \mid\ t_1]]
                     \] \[
                        \vdots
                     \] \[
                        \mathcal{M}[\Delta'; t_k] \colon [\Delta'] \to [\Delta', z_k \colon C_k[z_1 \mid\ t_1, \dots, z_{k-1} \mid\ t_{k-1}]]
                     \] are constructible by expansion. Then, by expansion rule \ref{rule:sort-symbol-intro-2},
                     $[\Delta] = [\Delta', y \colon S[z_1 \mid\ t_1, \dots, z_m \mid\ t_m]]$
                     is constructible by expansion.
                \item[] 
            \end{enumerate}
        \item[\ref{rule:cont-thin}] Since context $\Delta$ is of the form $\Gamma, y_1 \colon B_1, \dots, y_m \colon B_m$,
            then there exists an $k$, $1 \leq k \leq m$ such that $\Gamma, y_1 \colon B_1, \dots, y_{k-1} \colon B_{k-1}, y_{k+1} \colon B_{k+1}, \dots, y_m \colon B_m$
            is a context and $\Gamma, y_1 \colon B_1, \dots, y_{k-1} \colon B_{k-1} \vdash B_k \type$ is a derivable sequent.
            By induction hypothesis, if $\Delta' \equiv \Gamma, y_1 \colon B_1, \dots, y_{k-1}$, $[\Delta']$ is constructible by expansion,
            and by the same argument in the last point, so is $[\Delta', y_k \colon B_k]$.
            Finally, by considering the expansion rule \ref{rule:pullback-2} on $[\Delta', y_{k+1} \colon B_{k+1}, \dots, y_m \colon B_m]$
            and the projection morphism of $[\Delta', y_k \colon B_k]$ onto $[\Delta']$, we get that the pullback diagram
            \[
            \begin{tikzcd}
                {[\Delta]} 
                    \arrow[r, "{\langle y_1, \dots, y_{k-1}, y_{k+1}, \dots, y_m \rangle_\Gamma}"]
                    \arrow[d, open]
                & {[\Delta', y_{k+1} \colon B_{k+1}, \dots, y_m \colon B_m]}
                    \arrow[d, open]
                \\
                {[\Delta', y_k \colon B_k]} 
                    \arrow[r, open] 
                & {[\Delta']}
            \end{tikzcd}
            \] is constructible by expansion, hence, so is $[\Delta]$.
            % TODO : fix alignment
        \item[\ref{rule:cont-sub}] Then there exists a $k$, $1 \leq k \leq m$, a variable $x$ not appearing in $\Delta$, pre-type-expressions $A, B'_k, \dots, B'_m$ and a pre-object-expression $t$
            such that $\Gamma, y_1 \colon B_1, \dots, y_{k-1} \colon B_{k-1}, x \colon A, y_k \colon B'_k, \dots y_m \colon B'_m \ok$ and $\Gamma, y_1 \colon B_1, \dots, y_{k-1} \colon B_{k-1} \vdash t \colon A$
            are derivable sequents and $B_i = B'_i[x \mid\ t]$ for $k \leq i \leq m$.
            For simplicity of notation, we will note $\Delta' \equiv \Gamma, y_1 \colon B_1, \dots, y_{k-1} \colon B_{k-1}$
            Then, by induction hypothesis we have that both contexts $[\Delta', x \colon A]$ % TODO: check if this induction is valid
            and $[\Delta', x \colon A, y_k \colon B'_k, \dots y_m \colon B'_m]$
            and the section $[\langle y_1, \dots, y_{k-1}, t \rangle_\Gamma]$ are constructible by expansion.
            Hence, by rule \ref{rule:pullback-2}, so is the following pullback diagram:
            \[
            \begin{tikzcd}
                {[\Delta', y_k \colon B'_k[x \mid\ t], \dots, y_m \colon B'_m[x \mid\ t]]} 
                    \arrow[r, "{[\langle s, y_k, \dots, y_m \rangle_{\Delta'}]}"] 
                    \arrow[d, open]
                & {[\Delta', x \colon A, y_k \colon B'_k, \dots, y_m \colon B'_m]} 
                    \arrow[d, open]
                \\ 
                {[\Delta']}
                    \arrow[r, "{[\langle s \rangle_{\Delta'}]}"]
                & {[\Delta', x \colon A]}
            \end{tikzcd}
            \]
            Finally, $[\Delta] = [\Delta', y_k \colon B'_k[x \mid\ t], \dots, y_m \colon B'_m[x \mid\ t] ]$ 
            is constructible by expansion.

        \end{enumerate}

        Now consider a morphism \[
            [\langle t_1, \dots, t_w \rangle_\Gamma]
            \colon [\Gamma, y_1 \colon B_1, \dots, y_m \colon B_m]
            \to [\Gamma, z_1 \colon C_1, \dots, z_w \colon C_w]
        \] in $\mathcal{C}_{\Sigma, \Gamma}$. 
        We denote $\Delta \equiv \Gamma, y_1 \colon B_1, \dots, y_m \colon B_m$. 
        Then for $i$, $1 \leq i \leq w$, \[
            \Delta \vdash t_i \colon C_i[z_1 \mid\ t_1, \dots, z_{i-1} \mid\ t_{i-1}]
        \] are derivable sequents of $\text{CoC}_\Sigma$. Now we want to see that,
        for each $i$, the following section is constructible by expansion: \[
            [\langle y_1, \dots, y_m, t_i \rangle_\Gamma]
            \colon [\Delta] 
            \to [\Delta, z_i \colon C_i[z_1 \mid\ t_1, \dots, z_{i-1} \mid\ t_{i-1}]].
        \]
        We consider the ways the mentioned sequents are derived. 
        They must be derived by one of the following rules:
        \ref{rule:ass}, \ref{rule:refl-obj}, \ref{rule:conv1}, \ref{rule:prod-intro}, \ref{rule:prod-elim}, \ref{rule:prop-type} and \ref{rule:forall-intro}. 
        The cases of \ref{rule:refl-obj} and \ref{rule:conv1} are trivial. Hence, we consider the other cases:
        \begin{enumerate}
            \item[\ref{rule:ass}] There exists an $1 \leq j \leq m$ or $1 \leq k \leq n$ such that 
                \begin{align*}
                    t_i = y_j \quad \text{ orange } \quad t_i = x_k.
                \end{align*}
                We suppose that $t_i = y_j$ for some $j$, and the other case is equivalent.
                In that case $B_j = C_i [z_1 \mid\ t_1, \dots, z_{i-1} \mid\ t_{i-1}]$.
                We know that $\Delta$ is a context, hence by inductive hypothesis we have that $[\Delta]$ is constructible by expansion.
                When applying the expansion rule \ref{rule:variable-duplication-2}, we obtain that the 
                section \[
                    [\langle y_1, \dots, y_m, y_j \rangle_\Gamma] 
                    \colon [\Delta] 
                    \to [\Delta, y \colon B_j] 
                \] is constructible by expansion, and it is equal to $[\langle y_1, \dots, y_m, t_i \rangle_\Gamma]$,
                hence the least is also constructible by expansion.
            \item[\ref{rule:prod-intro}] Then there exists one pre-object-expression $t$ and two pre-type-expression $A, B$
                such that $t_i = (\lambda x \colon A) t$ and $C_i[z_1 \mid\ t_1, \dots z_{i-1} \mid\ t_{i-1}] = (\prod x \colon A) B$,
                for some variable $x$. Moreover, as $\level[\Delta; (\lambda x \colon A) t] = \level[\Delta, x \colon A ; t] + 1$,
                by induction hypothesis, the section \[
                    [\langle y_1, \dots, y_m, x, t\rangle_\Gamma]
                    \colon [ \Delta, x \colon A]
                    \to [ \Delta, x \colon A, y \colon B]
                \] is constructible by expansion.
                Applying the expansion rule \ref{rule:product-formation-2} to $[\Delta, x \colon A, y \colon B]$,
                we get that $[\Delta, f \colon (\prod x \colon A) B]$ is constructible by expansion.
                Finally, by the curry completion condition, we get that \[
                    \curry(\langle y_1, \dots, y_m, x, t \rangle_\Gamma) = \langle y_1, \dots, y_m, (\lambda x \colon A) t \rangle_\Gamma
                \] is constructible by expansion.
            \item[\ref{rule:prod-elim}] Then there exists one variable $x$, two pre-type-expressions $A, B$ and two pre-term-expressions $t, s$ such that
                $\Delta' \vdash t \colon (\prod x \colon A) B$, $\Delta' \vdash s \colon A$, $\Delta', x \colon A \vdash B \type$ and $B_m = \App([x \colon A] B, t, s)$.
                TODO: finish this, photo of partial proof 

            \item[\ref{rule:prop-type}] TODO
            \item[\ref{rule:forall-intro}] TODO  
        \end{enumerate}
\end{proof}